% 学术内容检测模块整体用例图（对应 03-5-content-detection.tex 全部功能点 FR-TPJC-0001～0003、FR-LWJC-0001～0005、FR-PLJC-0001～0002、FR-NRJC-0001～0002）
\begin{tikzpicture}[
  uc/.style={draw,ellipse,minimum width=2.2cm,minimum height=7mm,align=center,font=\tiny,inner sep=1pt},
  arr/.style={semithick}
]
  \node[draw,rounded corners,minimum width=11.1cm,minimum height=7.2cm] (frame) at (6.1,3.8) {};
  \node[anchor=north,yshift=-2pt] at (frame.north) {\small\textbf{学术内容检测模块（全部功能点）}};
  \node (actor) at (-0.3,3.8) {\shortstack{\small$\odot$\\[-1pt]\scriptsize 普通用户\\[-2pt]\tiny （提交检测任务）}};

  \node[uc] (u1) at (2.9,6.15) {提取图片\\FR-TPJC-0001};
  \node[uc] (u2) at (5.15,6.15) {学术图像AI鉴伪\\FR-TPJC-0002};
  \node[uc] (u3) at (7.4,6.15) {图像可疑区域标注\\FR-TPJC-0003};
  \node[uc] (u4) at (9.65,6.15) {全篇论文AIGC检测\\FR-LWJC-0001};

  \node[uc] (u5) at (2.9,4.75) {AI生成段落识别\\FR-LWJC-0002};
  \node[uc] (u6) at (5.15,4.75) {段落级AI生成概率\\FR-LWJC-0003};
  \node[uc] (u7) at (7.4,4.75) {事实性鉴伪子结论\\FR-LWJC-0004};
  \node[uc] (u8) at (9.65,4.75) {事实性鉴伪原因说明\\FR-LWJC-0005};

  \node[uc] (u9) at (2.9,3.35) {Review检测\\FR-PLJC-0001};
  \node[uc] (u10) at (5.15,3.35) {模板化评审倾向\\FR-PLJC-0002};
  \node[uc] (u11) at (7.4,3.35) {综合鉴伪结果生成\\FR-NRJC-0001};
  \node[uc] (u12) at (9.65,3.35) {检测结果可视化解释\\FR-NRJC-0002};

  \foreach \x in {u1,u2,u3,u4,u5,u6,u7,u8,u9,u10,u11,u12} {
    \draw[arr] (actor.east) -- ++(0.22,0) |- (\x.west);
  }
\end{tikzpicture}
